%%%%%%%%%%%%%%%%%%%%%%%%%%%%%%%%%%%%%%%%%%%%%%%%%%%%%%%%%%%%%%
% Artigo final de Trabalho de Conclusão de Curso             %
% Curso de Ciência da Computação -- Universidade de Passo Fundo
% Autor: Igor Zanette -- Orientador: Prof. Rafael Rieder     %
%                                                            %
% Classe CC-EC_TF-TC (modelo oficial do curso).              %
% Corpo do texto: 10 a 15 páginas, sem contar apêndices.     %
%%%%%%%%%%%%%%%%%%%%%%%%%%%%%%%%%%%%%%%%%%%%%%%%%%%%%%%%%%%%%%

\documentclass[alpha-refs,brazilian]{CC-EC_TF-TC}

\setlength {\marginparwidth }{2cm}
\AtBeginDocument{\renewcommand{\harvardand}{e}}

%%% Pacotes adicionais
\usepackage{lmodern}
\usepackage{float}
\usepackage{subcaption}
\usepackage{epstopdf}
\usepackage{todonotes}
\usepackage{multirow}
\usepackage{nicefrac}

%%% As figuras moram fora do artigo, em trabalho_de_conclusao/figuras/, porque são
%%% produzidas no Figma e reaproveitadas na apresentação. Exportar em PDF vetorial.
\graphicspath{{../figuras/}}

%%% Figura ainda não exportada: em vez de quebrar a compilação, desenha um quadro com o
%%% nome do arquivo esperado. Assim o texto pode ser revisado antes das figuras existirem.
%%% Aceita PDF (preferido, vetorial) ou PNG.
\newcommand{\figura}[2]{%
  \IfFileExists{../figuras/#1.pdf}%
    {\includegraphics[width=#2]{#1.pdf}}%
    {\IfFileExists{../figuras/#1.png}%
      {\includegraphics[width=#2]{#1.png}}%
      {\fbox{\parbox[c][3cm][c]{\dimexpr #2-2\fboxsep-2\fboxrule}{\centering\footnotesize
        \textcolor{red}{Figura ainda não exportada do Figma}\\[2pt]\texttt{figuras/#1.pdf}}}}}%
}

%%% Trecho que depende de dado ainda não coletado. Fica visível no PDF de propósito:
%%% um número de resultado nunca é escrito antes de existir.
\newcommand{\pendente}[1]{\textcolor{red}{[\textbf{PENDENTE:} #1]}}

%%% Hifenização de termos técnicos recorrentes
\hyphenation{
  apti-dão edá-fica frutí-feras espe-cia-lista con-cor-dân-cia infe-rência
  reco-men-dação inter-pre-tação Embra-pa fer-tili-dade adu-bação cala-gem
  rastre-abi-li-dade}

%%% Título provisório -- trocar aqui e no \titulores abaixo
\title{SIRAS: sistema especialista baseado em regras para recomendação de calagem e
adubação e estimativa de aptidão edáfica a partir da análise de solo}

\titleother{SIRAS: a rule-based expert system for liming and fertilization recommendation
and edaphic suitability estimation from soil analysis}

\author[1]{Igor Zanette}
\author[1]{Rafael Rieder}

\affil[1]{Universidade de Passo Fundo, Instituto de Tecnologia, Curso de Ciência da Computação}

\runningauthor{Zanette e Rieder, 2026}
\titulores{SIRAS: recomendação de calagem e adubação e aptidão edáfica}

\papercat{Trabalho de Conclusão de Curso}

\jyear{2026}
\jmonth{Dezembro}
\jcourse{Ciência da Computação}

\setcounter{page}{1}
\firstpage{1}

\begin{document}

\begin{frontmatter}

\maketitle
\thispagestyle{empty}

% Limite do modelo: 180 palavras cada; sem citações; 3 a 5 palavras-chave em ordem alfabética.

\begin{Abstract}
\textbf{Background:} turning a soil analysis into liming and fertilization recommendations
requires applying dozens of tables from the regional manual, and no tool available to technicians
in southern Brazil combines that recommendation with an assessment of the area's suitability. This
work developed SIRAS, a rule-based expert system that recommends liming and fertilization for 61
crops following the Rio Grande do Sul and Santa Catarina manual and estimates edaphic suitability
before and after correction, citing the rule and table behind every value.
\textbf{Results:} all 25 validation cases matched values calculated by hand from the manual (95\%
confidence interval: 86.7--100\%); the suitability module conformed to its frozen specification in
all 92 cases; and the system reproduced the 32 classifications published in the manual's annex.
Validation also revealed two defects missed by automated tests.
In a usability test, all eight participants completed a recommendation scenario unaided, and no
heuristic was rated above a cosmetic problem.
\textbf{Conclusions:} the system faithfully implements the manual's criteria and integrates
recommendation and suitability in a traceable report; the suitability criteria still require
independent agronomic validation.
\end{Abstract}

\begin{keywords}
Expert systems; liming and fertilization; rule-based systems; soil analysis; soil suitability
\end{keywords}

\begin{resumo}
\textbf{Contexto:} transformar a análise de solo em recomendação de calagem e adubação exige
aplicar dezenas de tabelas do manual regional, e nenhuma ferramenta acessível ao técnico gaúcho
integra essa recomendação à avaliação da aptidão da área. Este trabalho desenvolveu o SIRAS, um
sistema especialista baseado em regras que recomenda calagem e adubação para 61 culturas segundo o
Manual do RS e de SC e estima a aptidão edáfica antes e depois da correção, indicando a regra e a
tabela de cada valor.
\textbf{Resultados:} os 25 casos de validação concordaram com valores calculados à mão (IC 95\%:
86,7--100\%); a aptidão foi conforme à especificação nos 92 casos; e o sistema reproduziu as 32
classificações publicadas no anexo do Manual. A validação revelou dois defeitos que os testes automatizados não detectaram.
No teste de usabilidade, os oito participantes concluíram sozinhos um cenário de recomendação, e
nenhuma heurística recebeu nota acima de problema estético.
\textbf{Conclusões:} o sistema implementa fielmente os critérios do Manual e integra recomendação
e aptidão num laudo rastreável; os critérios de aptidão ainda requerem validação agronômica
independente.
\end{resumo}

\begin{palavras_chave}
Análise de solo; aptidão edáfica; calagem e adubação; sistemas baseados em regras; sistemas especialistas
\end{palavras_chave}


\end{frontmatter}
\clearpage

\section{Introdução}\label{sec:intro}

A interpretação da análise química e física do solo é a etapa que transforma um laudo de
laboratório em decisão de manejo: quanto calcário aplicar, quanto de nitrogênio, fósforo e
potássio fornecer e se a área está em condições de receber a cultura pretendida. O processo
exige conhecimento técnico e atenção rigorosa, porque erros de interpretação têm custo direto
para o produtor --- seja pelo uso desnecessário de insumos, seja pela perda de produtividade
decorrente de correções insuficientes. Sua complexidade advém da multiplicidade de variáveis
envolvidas, como pH, índice SMP, teores de fósforo e potássio, capacidade de troca de cátions
(CTC), saturação por bases e textura, e da necessidade de cruzá-las com as exigências de cada
cultura \citep{Donagema2011, ManualCalagem2016}.

No Rio Grande do Sul e em Santa Catarina, esses critérios estão consolidados no \textit{Manual
de Calagem e Adubação} da Comissão de Química e Fertilidade do Solo \citep{ManualCalagem2016},
que organiza a recomendação em dezenas de tabelas condicionadas ao grupo de cultura, ao sistema
de manejo, ao histórico da área e à expectativa de rendimento. Aplicá-lo à mão é trabalhoso e
sujeito a erro; e, na prática do crédito rural, observa-se que projetos agrícolas submetidos a
instituições financeiras podem apresentar recomendações inconsistentes, sem que haja um
instrumento objetivo e rastreável para conferi-las --- constatação que motivou este trabalho,
a partir da experiência do autor na assistência técnica a produtores e no atendimento bancário
a carteiras de crédito rural.

Há ferramentas que automatizam parte desse processo. O FertFácil \citep{Tiecher2016}
recomenda calagem e adubação pelos critérios do Manual e é amplamente usado no estado, e o
AdubaTec \citep{AdubaTec2021} faz o mesmo para mandioca e fruteiras tropicais, com recomendações da Embrapa. Em outra escala, o Sistema
de Avaliação da Aptidão Agrícola das Terras \citep{RamalhoFilhoBeek1995, EmbrapaAptidao2025}
classifica o potencial das terras em levantamentos regionais, mas não se presta à avaliação
de uma área individual a partir de sua análise de solo \citep{Hofig2015}. Existem, portanto,
sistemas que recomendam a correção do solo e métodos que avaliam a aptidão em escala regional,
mas não uma ferramenta acessível ao técnico gaúcho que integre as duas perspectivas a partir de
uma análise individual --- e que mostre, para cada número emitido, de que regra e de que
tabela ele veio.

Este trabalho teve por objetivo desenvolver e avaliar o \textbf{SIRAS} --- Sistema de
Interpretação e Recomendação Agronômica de Solo ---, um sistema especialista baseado em regras
que interpreta a análise de solo, recomenda calagem e adubação para 61 culturas e grupos de
culturas segundo o Manual \citep{ManualCalagem2016} e estima a aptidão edáfica da área para a
cultura escolhida, antes e depois da correção recomendada. A hipótese avaliada foi a de que o
SIRAS gera recomendações de calagem e adubação com taxa de concordância igual ou superior a
80\% em relação a valores de referência calculados de forma independente a partir do Manual.
A estimativa de aptidão, cujos critérios foram definidos pelo autor, foi verificada quanto à
conformidade da implementação com a especificação previamente congelada, e a interface foi
avaliada em teste com usuários do público-alvo.

O restante do artigo está organizado da seguinte forma: a \autoref{sec:revisao} apresenta os
fundamentos e os trabalhos relacionados; a \autoref{sec:metodos} descreve o método, os
instrumentos de avaliação e o sistema desenvolvido; a \autoref{sec:resultados} apresenta e
discute os resultados; e a \autoref{sec:conclusao} encerra o trabalho.

\section{Revisão de Literatura}\label{sec:revisao}

Esta seção reúne o necessário para compreender a solução: como a análise de solo é
interpretada, como o Manual de Calagem e Adubação transforma a interpretação em recomendação,
como se avalia a aptidão das terras e o que caracteriza um sistema especialista baseado em
regras. Ao final, posiciona-se o trabalho frente às ferramentas existentes.

\subsection{Análise de solo e sua interpretação}\label{sec:rev_analise}

A análise de solo é o principal instrumento de diagnóstico da fertilidade. Nela se determinam
o pH em água, que condiciona a disponibilidade da maioria dos nutrientes; os teores de fósforo
(P) e potássio (K) disponíveis; os teores trocáveis de cálcio (Ca), magnésio (Mg) e alumínio
(Al), este associado à toxidez em solos ácidos; a CTC, que expressa a capacidade do solo de
reter cátions; a saturação por bases (V\%); o teor de matéria orgânica (MO), relacionado ao
suprimento de nitrogênio; e o teor de argila, principal atributo físico considerado na
recomendação \citep{Donagema2011, ManualCalagem2016}.

Os valores, isoladamente, não dizem nada: a interpretação os compara a faixas de referência
que dependem do tipo de solo e da cultura \citep{EmbrapaDoc206}. No Manual RS/SC, os teores de
P e K são enquadrados em cinco classes de disponibilidade --- muito baixo, baixo, médio, alto
e muito alto ---, e cada classe corresponde a uma faixa de rendimento relativo da cultura; o
limite superior da classe médio é o teor crítico, acima do qual se obtém rendimento próximo
ao máximo econômico \citep{ManualCalagem2016}. As faixas não são universais: a interpretação
do fósforo depende da classe de teor de argila do solo, e a do potássio, da CTC a pH 7,0. Um
mesmo teor de P pode, assim, ser baixo em um solo argiloso e alto em um arenoso.

\subsection{Calagem e adubação no RS e em SC}\label{sec:rev_manual}

O Manual de Calagem e Adubação para os estados do Rio Grande do Sul e de Santa Catarina
\citep{ManualCalagem2016} é a referência técnica para correção e fertilização do solo na
região. A necessidade de calcário é estimada preferencialmente pelo \textbf{índice SMP}, que
indica a dose de calcário com PRNT 100\% necessária para elevar o pH da camada de 0 a 20~cm a
5,5, 6,0 ou 6,5; o método da saturação por bases é previsto como alternativa. Dois pontos
tornam essa recomendação menos direta do que uma consulta a tabela. Primeiro, o pH que dispara
a calagem e o pH alvo da dose não são dados pelo pH de referência da cultura, e frequentemente
divergem dele: são definidos por grupo de culturas e por sistema de manejo --- grãos em sistema convencional, por exemplo, só recebem
calcário quando o pH é inferior a 5,5, com dose para pH 6,0. Segundo, a dose é ajustada ao
contexto: em plantio direto consolidado aplica-se um quarto da dose, em superfície e limitada
a 5~t/ha; solos de baixo poder tampão usam equações com MO e Al em lugar da tabela; e a dose
é corrigida pelo PRNT do corretivo efetivamente disponível \citep{ManualCalagem2016}.

A adubação fosfatada e potássica é dimensionada pela classe de disponibilidade e, para várias
culturas, pela expectativa de rendimento; a nitrogenada, pelo teor de MO e pela cultura
antecedente, e é dispensada nas leguminosas, que fixam nitrogênio biologicamente. O Manual
publica tabelas próprias para cada cultura ou grupo, com estruturas distintas: grãos separam
adubação de correção e de manutenção, frutíferas distinguem pré-plantio, crescimento e
produção, e a erva-mate e o tabaco exigem variáveis específicas de manejo e de tipo
\citep{ManualCalagem2016}.

\subsection{Aptidão das terras e o recorte edáfico}\label{sec:rev_aptidao}

O Sistema de Avaliação da Aptidão Agrícola das Terras \citep{RamalhoFilhoBeek1995} classifica
as terras pelo grau de limitação --- nulo, ligeiro, moderado, forte e muito forte --- que cada
fator impõe, considerando deficiência de fertilidade, deficiência hídrica, excesso de água,
suscetibilidade à erosão e impedimentos à mecanização. A classe resulta do fator mais
limitante, avaliado nas condições atuais e após a aplicação de práticas de melhoria
\citep{Hofig2015}. É o sistema empregado no mapeamento nacional \citep{EmbrapaAptidao2025}.
Sua escala, porém, é regional: \cite{Hofig2015} observam que ele não indica práticas de manejo
e requer adaptações para o planejamento de propriedades individuais.

Uma alternativa é restringir a avaliação aos atributos que a própria análise de solo informa,
cruzados com as exigências da cultura, sem relevo e sem clima. Neste trabalho, essa estimativa
é chamada de \textbf{aptidão edáfica}, para que não seja confundida com a aptidão agrícola das
terras: ela herda a arquitetura de graus de limitação e de fator mais limitante, mas não
reproduz o sistema oficial nem equivale a ele.

\subsection{Sistemas especialistas baseados em regras}\label{sec:rev_se}

Sistemas especialistas reproduzem o raciocínio de um especialista humano em um domínio
restrito \citep{Mendes1997}. São compostos por uma \textbf{base de conhecimento}, que
armazena fatos e regras do domínio; um \textbf{motor de inferência}, que aplica as regras aos
dados de entrada para derivar conclusões; e uma \textbf{interface} com o usuário. Na forma mais
difundida, o conhecimento é representado por regras SE-ENTÃO, e o motor percorre a base para
identificar as condições satisfeitas pelos dados \citep{Lanzer2009}. A representação explícita
das regras é o que permite ao sistema \emph{explicar} uma conclusão --- característica central
em domínios nos quais a recomendação precisa ser justificada a terceiros.

Na agricultura, \cite{Lanzer2009} aplicou representação hierárquica do conhecimento à
classificação do uso e manejo das terras e observou que, quando os estados do domínio já estão
definidos em sistemas de classificação publicados --- como as tabelas do Manual ---, a
construção da base é facilitada. \cite{Shahzadi2016}, em outro domínio, alimentaram regras
especialistas com dados de sensores para recomendar o manejo de pragas e doenças, evidenciando a
aplicabilidade da abordagem ao suporte à decisão agronômica.

\subsection{Trabalhos relacionados}\label{sec:rev_relacionados}

O \textbf{FertFácil} \citep{Tiecher2016} é um sistema on-line e gratuito que recomenda calagem
e adubação pelos critérios do Manual para mais de cem cultivos do RS, SC, MS e Paraguai, adotado
pela Emater/RS; hoje incorpora também o manual do Paraná \citep{FertFacil2026}. O \textbf{AdubaTec} \citep{AdubaTec2021}, da Embrapa Mandioca e Fruticultura e da
UFRB, recomenda calagem e adubação para mandioca e fruteiras tropicais a partir das recomendações
publicadas pela própria Embrapa, com arquitetura parametrizável que permite incluir novas culturas
sem alterar o código. O \textbf{Sistema de Avaliação da Aptidão
Agrícola das Terras} \citep{RamalhoFilhoBeek1995, EmbrapaAptidao2025} classifica o potencial
das terras em escala regional. A \autoref{tab:relacionados} sintetiza o posicionamento: os
sistemas existentes tratam separadamente a recomendação de manejo e a avaliação de aptidão, e
nenhum integra as duas a partir de uma análise individual de solo sob os critérios do RS e SC.

\begin{table}[htb]
\centering
\caption{Posicionamento do SIRAS frente aos trabalhos relacionados.}
\label{tab:relacionados}
\footnotesize
\begin{tabularx}{\columnwidth}{@{}lCCCC@{}}
\toprule
 & \textbf{Critérios} & \textbf{Escala} & \textbf{Calagem e adubação} & \textbf{Aptidão} \\
\midrule
FertFácil  & RS/SC e PR & área & sim & não \\
AdubaTec   & Embrapa & área  & sim & não \\
SAAAT      & nacional & regional & não & sim \\
\textbf{SIRAS} & \textbf{RS/SC} & \textbf{área} & \textbf{sim} & \textbf{sim (edáfica)} \\
\bottomrule
\end{tabularx}

\vspace{2pt}
\raggedright\scriptsize SAAAT: Sistema de Avaliação da Aptidão Agrícola das Terras.
\end{table}

\section{Materiais e Métodos}\label{sec:metodos}

Esta seção caracteriza a pesquisa, descreve a construção da base de conhecimento e dos
critérios de aptidão, os procedimentos de validação e de avaliação de usabilidade e, por fim,
o sistema desenvolvido.

\subsection{Caracterização da pesquisa}\label{sec:caracterizacao}

A pesquisa é aplicada e descritiva: estrutura um conjunto de critérios técnicos consolidados e
os aplica na construção e na avaliação de um artefato computacional \citep{Wazlawick2014}.
Quanto ao estilo, enquadra-se na apresentação de um produto, estilo que, segundo
\cite{Wazlawick2014}, exige comparação com alguma referência, pois a mera demonstração de
funcionamento tem baixa aceitação científica. Por isso, as saídas do SIRAS foram confrontadas
com um oráculo calculado de forma independente da implementação, e a interface foi submetida a
usuários do público-alvo. A validação das recomendações é quantitativa; a avaliação de
usabilidade, descritiva e qualitativa.

\subsection{Base de conhecimento e escopo}\label{sec:base}

A fonte única das recomendações é o Manual de Calagem e Adubação \citep{ManualCalagem2016}.
Suas tabelas de interpretação, de calagem e de adubação foram transcritas manualmente pelo autor
para 20 arquivos JSON, organizados em dados comuns --- critérios de calagem, tabela SMP,
interpretação de P e K, pH de referência, critérios de aptidão --- e em um diretório por grupo de
culturas. Cada arquivo registra a tabela e a página de origem e é verificado, a cada carga, por um
esquema JSON Schema próprio, de modo que um valor fora do domínio ou um campo ausente impede a
execução em vez de produzir uma recomendação silenciosamente errada.

O escopo compreende as \textbf{61 culturas e grupos de culturas} para os quais o Manual publica
tabelas completas de adubação, em seis grupos (\autoref{fig:escopo}): grãos (21), hortaliças
(18), tubérculos (2), frutíferas (17), erva-mate (1) e outras culturas comerciais --- cana-de-açúcar
e tabaco (2). A contagem segue a organização do Manual: culturas que compartilham uma tabela,
como melancia e melão, formam uma entrada. Ficaram fora do escopo de recomendação o arroz
irrigado, a mandioca e as espécies florestais, cuja lógica de calagem difere do modelo adotado,
e, por delimitação, forrageiras, plantas medicinais e ornamentais e sistemas especiais de
produção.

\begin{figure*}[htb]
\centering
\figura{F3-escopo}{\textwidth}
\caption{Escopo do SIRAS: os seis grupos de culturas, as particularidades de recomendação de
cada um e o que ficou fora do escopo.}
\label{fig:escopo}
\end{figure*}

A transcrição é o ponto mais sensível do trabalho, porque um número errado na base contamina
todas as recomendações que dependem dele. Além da conferência manual, os grupos de exigência
de P e K de cada cultura foram transcritos três vezes, por caminhos distintos e em momentos
distintos, e um teste automatizado exige que as três transcrições concordem.

\subsection{Critérios de aptidão edáfica}\label{sec:ccae}

O Manual não define classes de aptidão. Os critérios foram, portanto, construídos pelo autor e
reunidos em uma especificação própria, o \textbf{Caderno de Critérios de Aptidão Edáfica}
(CCAE), segundo três princípios: nenhum limiar numérico é criado pelo autor --- todos provêm das
tabelas de interpretação do Manual \citep{ManualCalagem2016} ou de \cite{EmbrapaDoc206}, com
tabela e página registradas; a arquitetura de graus de limitação e de fator mais limitante vem
do sistema oficial \citep{RamalhoFilhoBeek1995}, conforme descrito por \cite{Hofig2015}; e toda
escolha não determinada pelas fontes é isolada como decisão do autor, registrada e implementada
como parâmetro configurável.

Para três dessas decisões buscou-se respaldo na literatura. Os limites de saturação por alumínio
adotados coincidem com os de \cite{Alvarez1999}, e tratar Ca e Mg como corrigíveis pela própria
calagem é compatível com \cite{Sousa1989}. Para o teto de dose que define quando a correção da
acidez é exequível em uma safra não se encontrou equivalente: a literatura trata do
parcelamento, e não de um limite máximo. Esse resultado negativo é declarado, e a decisão
permanece do autor.

O CCAE foi versionado: cada versão foi aprovada pelo orientador e congelada com marcação no
repositório, e alterações posteriores geraram versão nova, com o motivo registrado. O gabarito
dos casos de teste foi produzido aplicando a versão vigente à mão, sem consultar o código; só
assim a verificação mede se a implementação reproduz a especificação, e não se o código
concorda consigo mesmo.

\subsection{Validação das recomendações}\label{sec:validacao}

A hipótese do trabalho foi avaliada pela \textbf{concordância} entre as recomendações do SIRAS e
valores de referência. Um caso é concordante quando \emph{todas} as saídas declaradas na
referência --- necessidade de calcário (t/ha) e doses de N, P\textsubscript{2}O\textsubscript{5}
e K\textsubscript{2}O (kg/ha) --- coincidem com as do sistema, admitida apenas a tolerância de
arredondamento: calcário com uma casa decimal e nutrientes em quilogramas inteiros. Acertar três
de quatro saídas é discordância.

Os valores de referência constituem o \textbf{oráculo} e foram calculados pelo autor, à mão,
aplicando diretamente as tabelas do Manual, sem executar o sistema
(\autoref{fig:oraculo}). Sem essa independência, a validação mediria apenas a consistência do
código consigo mesmo. O conjunto final reuniu 25 casos distribuídos pelos seis grupos,
escolhidos para cobrir as classes de interpretação e os casos especiais --- plantio direto
consolidado, baixo poder tampão, fases das frutíferas, erva-mate e tipos de tabaco. Como o
sistema é determinístico, a análise consiste na proporção de casos concordantes, com intervalo
de confiança de 95\% pelo método de \cite{Wilson1927}, adequado a proporções próximas de 100\% e
a amostras pequenas. A hipótese seria aceita com taxa igual ou superior a 80\%.

\begin{figure}[htb]
\centering
\figura{F13-oraculo}{\columnwidth}
\caption{Independência do oráculo: o Manual alimenta dois caminhos separados --- o cálculo à mão
e a implementação ---, que só se encontram na comparação.}
\label{fig:oraculo}
\end{figure}

A aptidão edáfica não integra a hipótese: como seus critérios foram definidos pelo autor, medir
concordância com eles seria circular. Ela foi submetida a \textbf{verificação de conformidade}
--- se a implementação reproduz o CCAE congelado --- sobre 92 casos com gabarito produzido à mão,
dos quais 88 com classe definida e 4 que devem resultar em estado indeterminado por dado
ausente. Doze desses casos usam as glebas reais publicadas no Anexo~1 do Manual, cujas classes de
P e K o próprio Manual divulga e servem, assim, de oráculo externo para a etapa de interpretação.

O processo de desenvolvimento foi apoiado por testes automatizados (5.404 ao final), que
incluem os casos de validação; por um registro de decisões de arquitetura (13 decisões
documentadas, com contexto e alternativas); e por versionamento em Git.

\subsection{Avaliação de usabilidade}\label{sec:usabilidade}

A interface foi avaliada em teste com 5 a 8 usuários do público-alvo --- agrônomos, técnicos
agrícolas e estudantes concluintes de agronomia com contato com análise de solo ---, precedido
de um piloto fora da amostra. O tamanho segue \cite{NielsenLandauer1993}, segundo os quais cerca
de cinco participantes revelam a maior parte dos problemas de uma interface. Cada sessão, de
cerca de 35 minutos, percorreu três roteiros: uma demonstração conduzida pelo pesquisador sobre
uma análise real do Anexo~1 do Manual; um cenário fechado, entregue em ficha impressa e
executado pelo participante sozinho, cujas três perguntas têm resposta conhecida; e o uso livre,
com cultura e dados escolhidos pelo participante (\autoref{fig:sessao}).

Depois da sessão, sem a presença do pesquisador, o participante respondeu a um questionário
estruturado pelas dez heurísticas de \cite{Nielsen1994}, atribuindo a cada uma a gravidade dos
problemas encontrados, de 0 (não é problema) a 4 (catástrofe), ou indicando que não passou pela
situação, e descrevendo as falhas e a tela em que ocorreram. Quatro perguntas abertas encerram
o instrumento e separam o juízo sobre a ferramenta do juízo sobre a recomendação agronômica.

O instrumento substituiu o \textit{System Usability Scale} \citep{Brooke1996} previsto na proposta,
porque o SUS produz um escore único, e esta etapa precisava de um diagnóstico que localizasse os
problemas e os ordenasse por gravidade. Não se trata de avaliação heurística no sentido de
\cite{NielsenMolich1990}, que é inspeção por especialistas sem usuários: as notas são atribuídas
por usuários a heurísticas e exprimem percepção de gravidade, analisada por estatística
descritiva. A versão avaliada foi congelada antes da primeira sessão e mantida até a última.

\subsection{O sistema desenvolvido}\label{sec:siras}

Do ponto de vista do usuário, o SIRAS tem três telas: escolher a cultura, informar a análise e
ler o laudo. A simplicidade é deliberada e esconde a maior parte do sistema
(\autoref{fig:iceberg}): cada laudo resulta da resolução de critérios por grupo, de dezenas de
tabelas transcritas, de dois cenários de aptidão e de um registro de cada decisão tomada.

\begin{figure*}[htb]
\centering
\figura{F1-iceberg}{\textwidth}
\caption{O que o usuário vê e o que o sistema faz: as três telas da interface sobre as camadas
de conhecimento, inferência e verificação que produzem o laudo.}
\label{fig:iceberg}
\end{figure*}

\subsubsection{Arquitetura e tecnologias}

O sistema foi implementado em Python~3.12 com o microframework Flask~3 para a camada
web\footnote{Disponíveis em \url{https://www.python.org} e
\url{https://flask.palletsprojects.com}.}, com HTML, CSS e JavaScript na interface e as
bibliotecas \texttt{jsonschema} e \texttt{pytest} para validação e testes. Não há banco de dados,
autenticação nem chamada de rede: a aplicação roda localmente, no navegador, e é iniciada por um
único comando que prepara o ambiente e sobe o servidor.

A arquitetura segue a separação clássica dos sistemas especialistas
\citep{SantosEtAl1994, Lanzer2010}. A \textbf{base de conhecimento} são os arquivos JSON; o
\textbf{motor de inferência} é um pacote em Python puro que não depende do Flask e pode ser
executado e testado isoladamente; e a \textbf{interface} é uma camada fina que apenas coleta os
dados e apresenta o laudo. O motor foi escrito sem bibliotecas de inferência, como Experta ou
PyKE, para que cada regra permanecesse legível e auditável. Todo o motor é acessado por uma
única função, pura e determinística, que recebe a análise, a cultura e o contexto de manejo e
devolve o laudo. As camadas e suas dependências estão na \autoref{fig:arquitetura}, no
apêndice.

\subsubsection{Fluxo de inferência e rastreabilidade}

A \autoref{fig:pipeline} apresenta o percurso de uma análise pelo motor. Os dados são validados
quanto ao domínio de cada variável; os teores de P e K são classificados pela classe de argila e
pela CTC (\autoref{fig:interpretacao}, no apêndice); a calagem, a adubação e a formulação do adubo são calculadas; e a aptidão é avaliada
em dois cenários. Cada passo registra, em uma \textbf{trilha de inferência}, a regra aplicada,
as entradas, a saída e a fonte no Manual, com tabela e página. É essa trilha que o laudo
apresenta ao usuário na seção ``de onde vem cada número'' (\autoref{fig:trilha}), e é ela que
permitiu localizar a causa de cada divergência durante a validação.

\begin{figure*}[htb]
\centering
\figura{F4-pipeline}{\textwidth}
\caption{Percurso de uma análise de solo pelo motor de inferência. Cada etapa registra na
trilha a regra aplicada, as entradas, a saída e a fonte no Manual.}
\label{fig:pipeline}
\end{figure*}

\subsubsection{Calagem}

A calagem não é uma consulta direta à tabela SMP (\autoref{fig:calagem}). A cultura, o sistema
de manejo e a condição da área determinam primeiro um \textbf{critério}, entre os definidos nas
Tabelas 5.3 a 5.7 do Manual, que fixa a condição de disparo, o pH alvo da dose e o fator
aplicado. Só então o índice SMP é consultado (Tabela 5.2), e a dose recebe os ajustes: fator do
manejo, correção pelo PRNT do corretivo ($\mathit{dose} = \mathit{NC} \times 100 /
\mathit{PRNT}$), profundidade e limite de aplicação superficial. Os solos de baixo poder tampão
seguem as equações do Manual em lugar da tabela, e a erva-mate, sem pH de referência, segue o
critério de saturação por bases. Fora da faixa tabelada do índice SMP, o sistema recusa a
entrada em vez de extrapolar.

\begin{figure}[htb]
\centering
\figura{F5-calagem}{\columnwidth}
\caption{Árvore de decisão da calagem, com o caso da Gleba~4 do Anexo~1 do Manual em plantio
direto consolidado: 6,8~t/ha pela Tabela~5.2, 1,7~t/ha pelo fator de um quarto e 2,3~t/ha com
calcário de PRNT~75\%.}
\label{fig:calagem}
\end{figure}

\subsubsection{Adubação e formulação do adubo}

As doses de N, P\textsubscript{2}O\textsubscript{5} e K\textsubscript{2}O seguem dois caminhos.
Os grãos têm algoritmo próprio, que combina adubação de correção e de manutenção, o número do
cultivo, a expectativa de rendimento e, para o nitrogênio, o teor de MO e a cultura antecedente.
Os demais grupos recebem a dose publicada por classe de teor, parametrizada pelas variáveis que
cada um exige: a fase das frutíferas, o programa e o manejo da erva-mate, o tipo de tabaco e o
ciclo da cana (\autoref{fig:adubacao}).

Atendendo a pedido do orientador, o laudo indica também \emph{qual adubo comprar}. Como não há
catálogo oficial das formulações comercializadas, o sistema apresenta três respostas
complementares: a \textbf{fórmula ideal}, na proporção exata da recomendação, pelo método do
Manual; as \textbf{melhores fórmulas de um catálogo} de 38 formulações comerciais, cada uma com
sua dose e os complementos necessários; e a \textbf{mistura de adubos simples}, sempre exata.
Para cada fórmula do catálogo, a dose é escolhida entre os pontos em que ela atende exatamente
o P\textsubscript{2}O\textsubscript{5} ou o K\textsubscript{2}O recomendado; como a função que
penaliza excesso, falta e operações de complemento é convexa por partes, seu mínimo está
necessariamente em um desses pontos, e um deles é a regra do Manual. O resultado, portanto,
nunca é pior que o do Manual sob o critério adotado. Os pesos dessa função são decisão do
autor e estão declarados como tal. A \autoref{fig:formulacao} ilustra os três caminhos e a
escolha da dose.

\subsubsection{Aptidão edáfica}

A aptidão é avaliada por sete fatores de limitação --- acidez, fósforo, potássio, cálcio e
magnésio, CTC, matéria orgânica e textura ---, cada um com grau de nulo a muito forte
(\autoref{fig:aptidao}). A classe resulta do fator mais limitante, como no sistema oficial:
apta, apta com restrições, restrita ou inapta sem correção prévia. O caderno prevê ainda uma
regra de rebaixamento quando três ou mais fatores são ao menos moderados, mas, por não ter
respaldo nas fontes, ela foi implementada desativada e serve apenas à análise de sensibilidade.
Dado obrigatório ausente nunca é imputado: produz o estado indeterminado.

A avaliação é feita em dois cenários: o \textbf{atual}, sobre a análise como está, e o
\textbf{potencial}, após a aplicação integral das recomendações emitidas pelo próprio SIRAS.
Fatores corrigíveis, como acidez, fósforo e potássio, deixam de limitar no cenário potencial se a
correção for exequível; fatores permanentes, como CTC e textura, não mudam. A diferença entre os
dois cenários mede o ganho atribuível à recomendação, e é nela que os dois módulos do sistema se
integram.

\begin{figure*}[htb]
\centering
\figura{F9-aptidao}{\textwidth}
\caption{Estimativa de aptidão edáfica: sete fatores, graus de limitação, fator mais limitante e
os cenários atual e potencial, ilustrados pela Gleba~4 do Anexo~1 do Manual com milho (classe
restrita, determinada pelo potássio).}
\label{fig:aptidao}
\end{figure*}

\subsubsection{Interface e laudo}

O fluxo de uso está na \autoref{fig:fluxo}, no apêndice. Na escolha da cultura, as 61 opções aparecem por grupo. No
formulário, a classe de cada teor é exibida enquanto o usuário digita, os campos condicionais
surgem conforme a cultura escolhida, e várias áreas podem ser informadas em um mesmo laudo,
compartilhando cultura e manejo. Enquanto o laudo é gerado, uma tela de espera mostra as etapas
que o motor percorre.

O laudo (\autoref{fig:laudo}) foi organizado em duas partes depois que uma medição mostrou que,
na primeira versão, 71\% do texto de um laudo de uma área eram complementos, e que um laudo de seis
áreas ocupava 27 folhas. A parte essencial traz o quadro comparativo das áreas, as quatro doses,
a interpretação dos atributos, a aptidão e o total a comprar; o anexo traz a formulação do adubo,
a forma de aplicação e a trilha de inferência, com cada regra listada uma única vez. O documento
pode ser impresso ou salvo em PDF pelo navegador, e a análise pode ser editada e o laudo gerado
de novo sem redigitação.

\begin{figure*}[htb]
\centering
\figura{F11-laudo}{\textwidth}
\caption{Estrutura do laudo emitido pelo SIRAS: parte essencial e anexo.}
\label{fig:laudo}
\end{figure*}

\section{Resultados e Discussão}\label{sec:resultados}

Esta seção apresenta, para cada dimensão avaliada, o resultado obtido e a sua interpretação: a
concordância das recomendações, os defeitos que a validação revelou, a conformidade da aptidão
edáfica e sua sensibilidade às decisões do autor, a avaliação de usabilidade e as limitações.

\subsection{Concordância das recomendações}\label{sec:res_concordancia}

Os 25 casos de validação foram concordantes, com taxa de 100\% e intervalo de confiança de 95\%
de [86,7\%; 100,0\%] (\autoref{tab:concordancia}). Como o limite inferior do intervalo supera os
80\% fixados na hipótese, a hipótese foi aceita. O dado mais importante, porém, não é a taxa, e
sim o intervalo: com os 15 casos que o conjunto tinha até 10/10/2026, a taxa já era de 100\%,
mas o intervalo era de [79,6\%; 100,0\%], com o limite inferior abaixo da meta. O problema não
era de acerto, e sim de tamanho de amostra, e foi resolvido com dez casos novos.

\begin{table}[htb]
\centering
\caption{Concordância das recomendações de calagem e adubação por grupo de culturas.}
\label{tab:concordancia}
\footnotesize
\begin{tabular}{lccc}
\toprule
\textbf{Grupo} & \textbf{Casos} & \textbf{Concord.} & \textbf{IC 95\% (Wilson)} \\
\midrule
Grãos          & 12 & 12 & [75,8\%; 100\%] \\
Frutíferas     &  4 &  4 & [51,0\%; 100\%] \\
Hortaliças     &  3 &  3 & [43,9\%; 100\%] \\
Cana e tabaco  &  3 &  3 & [43,9\%; 100\%] \\
Erva-mate      &  2 &  2 & [34,2\%; 100\%] \\
Tubérculos     &  1 &  1 & [20,7\%; 100\%] \\
\midrule
\textbf{Total} & \textbf{25} & \textbf{25} & \textbf{[86,7\%; 100\%]} \\
\bottomrule
\end{tabular}
\end{table}

Dois pontos limitam a leitura desse resultado. Primeiro, o conjunto ficou abaixo dos 40 casos
previstos na proposta: cada caso exige o cálculo à mão de todas as saídas, com memória
registrada, e os casos foram escolhidos por lacuna de cobertura, e não por quantidade. Os
intervalos por grupo, largos, mostram o que isso custa: a evidência é forte para o conjunto e
para os grãos, e fraca para os grupos com um a quatro casos. Segundo, sistema e oráculo partem do
mesmo Manual. A concordância atesta que a transcrição e a implementação reproduzem os critérios
publicados, e não que esses critérios sejam agronomicamente adequados a uma área específica.

O Anexo~1 do Manual oferece uma verificação que não depende do autor: ele publica a análise de
quatro glebas e as classes de P e K resultantes para alho, milho, trigo e eucalipto. O motor do
SIRAS reproduziu as 32 classificações publicadas, sem divergência, e a tabela SMP transcrita
reproduziu as três doses de calcário que o texto do anexo calcula.

\subsection{Defeitos revelados pela validação}\label{sec:res_defeitos}

Os dez casos acrescentados para estreitar o intervalo cobriram grupos e situações que o conjunto
inicial não exercitava --- dos 15 casos originais, nove eram de calagem e um só passava pelo
algoritmo de adubação dos grãos. Dois deles divergiram do sistema, e nos dois a falha era do
sistema, não do caso. Na adubação nitrogenada, o adicional de N por rendimento esperado estava
transcrito na base para todas as culturas a que se aplica, mas nenhuma linha de código o lia: a
dose de N não variava com a expectativa de rendimento, enquanto a de P e K variava. Na classe
\emph{muito alto} de P e K, no segundo cultivo, o sistema aplicava a dose de manutenção onde o
Manual indica a de reposição, porque o texto transcrito juntava as duas numa mesma frase.

Os dois defeitos estavam no sistema desde a implementação e haviam passado por mais de cinco mil
testes automatizados sem que nenhum falhasse, porque nenhum teste combinava as entradas que os
revelam. É o resultado mais instrutivo da validação: uma taxa de 100\% sobre um conjunto
concentrado no caminho mais exercitado não media o que parecia medir. O valor do oráculo
calculado à mão está em que ele é escrito a partir do Manual, e não do código, e por isso pode
discordar do código.

\subsection{Conformidade da aptidão edáfica}\label{sec:res_aptidao}

A implementação reproduziu o CCAE nos 88 casos com classe definida, tanto na classe quanto no
fator determinante, e os quatro casos de dado ausente resultaram, como especificado, no estado
indeterminado. O resultado atesta a correção e a rastreabilidade da implementação, e não a
adequação agronômica dos critérios, que não foram validados por especialistas independentes nem em
campo.

Para medir quanto a classe depende das decisões do autor, o conjunto foi executado em quatro
configurações (\autoref{tab:sensibilidade}). A regra de rebaixamento por acúmulo, a decisão de
maior fragilidade do caderno, é também a de maior efeito: ativada, mudaria 12 dos 88 casos, quase
todos de \emph{restrita} para \emph{inapta sem correção prévia}, e dobraria essa última classe.
Todas as mudanças ocorreriam no cenário atual, nunca no potencial, porque após a correção restam
apenas os fatores permanentes. Esse efeito justifica mantê-la desativada no resultado oficial.
Os fatores permanentes pesam pouco: retirar a textura, ou a textura e a CTC, muda dois e cinco
casos, todos de \emph{apta com restrições} para \emph{apta}.

\begin{table}[htb]
\centering
\caption{Distribuição das classes de aptidão nos 88 casos, por configuração.}
\label{tab:sensibilidade}
\footnotesize
\begin{tabular}{lccccc}
\toprule
\textbf{Configuração} & \textbf{A} & \textbf{AR} & \textbf{R} & \textbf{I} & \textbf{Mudam} \\
\midrule
Sistema (fator mais limitante) & 29 & 30 & 18 & 11 & --- \\
Com rebaixamento por acúmulo   & 29 & 29 &  8 & 22 & 12 \\
Sem textura (F7)               & 31 & 28 & 18 & 11 &  2 \\
Sem CTC e textura (F5 e F7)    & 34 & 25 & 18 & 11 &  5 \\
\bottomrule
\end{tabular}

\vspace{2pt}
\raggedright\scriptsize A: apta; AR: apta com restrições; R: restrita; I: inapta sem correção prévia.
\end{table}

A comparação entre cenários é a parte da aptidão que integra os dois módulos do sistema. No caso
da Gleba~4 com milho (\autoref{fig:aptidao}), a área é \emph{restrita} hoje, limitada pelo
potássio, e passa a \emph{apta} após a correção recomendada pelo próprio SIRAS, quando a única
limitação remanescente é a CTC, que não se corrige com insumos. O laudo deixa assim de dizer
apenas quanto aplicar e passa a dizer o que a aplicação muda.

\subsection{Avaliação de usabilidade}\label{sec:res_usabilidade}
Participaram oito pessoas: cinco técnicos agrícolas, dois estudantes de agronomia e um
engenheiro agrônomo. Quatro têm mais de dez anos de trabalho com análise de solo ou recomendação
de adubação; os demais se distribuem entre seis e dez anos (um), dois a cinco anos (um), menos de
dois anos (um) e ainda sem atuação profissional (um). Hoje, cinco fazem a recomendação à mão, com
o Manual, quatro usam o sistema do laboratório, dois um programa de cooperativa ou revenda e um
uma planilha própria --- a pergunta admitia mais de uma resposta. Todos se declararam
confortáveis (quatro) ou muito confortáveis (quatro) com ferramentas digitais, nenhum conhecia o
SIRAS antes do teste e um deles (P7) já tinha usado o FertFácil.

Os oito concluíram sozinhos o cenário fechado do segundo roteiro, sem intervenção do pesquisador.
No questionário, das 80 avaliações possíveis (oito participantes, dez heurísticas), 79 receberam
nota; em uma o participante indicou não ter passado pela situação. Foram 73 notas 0 (\emph{não é
problema}) e seis notas 1 (\emph{problema estético}); nenhuma heurística recebeu nota 2 ou maior
(\autoref{tab:heuristicas}). As seis notas 1 vieram de três participantes, e quatro delas de um só
(P5). A mediana foi 0 em todas as heurísticas, e a ordenação por gravidade que o protocolo previa
não se aplica: as diferenças entre heurísticas são de uma nota, de um participante.

\begin{table}[htb]
\centering
\caption{Gravidade atribuída pelos participantes a cada heurística ($n=8$).}
\label{tab:heuristicas}
\footnotesize
\begin{tabular}{llccc}
\toprule
\# & Heurística & Nota 0 & Nota 1 & Não observou \\
\midrule
H1  & Visibilidade do status do sistema      & 7 & 1 & -- \\
H2  & Correspondência com o mundo real       & 7 & 1 & -- \\
H3  & Controle e liberdade do usuário        & 8 & -- & -- \\
H4  & Consistência e padrões                 & 8 & -- & -- \\
H5  & Prevenção de erros                     & 7 & 1 & -- \\
H6  & Reconhecimento em vez de memorização   & 8 & -- & -- \\
H7  & Flexibilidade e eficiência de uso      & 8 & -- & -- \\
H8  & Estética e design minimalista          & 7 & 1 & -- \\
H9  & Ajuda para reconhecer e corrigir erros & 6 & 1 & 1 \\
H10 & Ajuda e documentação                   & 7 & 1 & -- \\
\bottomrule
\end{tabular}

\vspace{2pt}
\raggedright\scriptsize Escala de 0 (não é problema) a 4 (catástrofe). Nenhuma heurística
recebeu nota 2, 3 ou 4.
\end{table}

Os relatos em texto são mais informativos que as notas. Dois descrevem a mesma falha, de
prevenção e recuperação de erro no formulário: P5 observou que a tecla Enter gera o laudo mesmo
com dados incompletos e sugeriu uma confirmação; P7, que de início não informou um campo e recebeu
apenas uma mensagem de erro, sem indicação do que faltava, e pediu que os campos pendentes fossem
destacados. Os dois pontos foram corrigidos depois da coleta, fora da versão avaliada: a tecla
Enter passou a levar ao campo seguinte, e a recusa de um valor passou a nomear e destacar o campo.
P2 relatou ter conseguido informar um só tipo de adubo da revenda, embora o formulário aceite até
três fórmulas. Dois participantes apontaram o contexto de uso: P1 lembrou que o teste não reproduz
o dia a dia com internet lenta, e P5 pediu uma versão que funcione sem conexão e fora do
computador do escritório.

Nas perguntas abertas, os oito disseram que usariam a ferramenta no trabalho, com justificativas
centradas em tempo e precisão --- \emph{``ter o laudo completo em menos de 10 min''} (P1),
\emph{``evita erros, e possui linguagem agronômica de fácil entendimento''} (P4). O participante
que já conhecia o FertFácil destacou justamente o que o diferencia: \emph{``as informações de onde
saiu os números achei muito bom''} (P7). Sobre a recomendação, os oito a consideraram confiável e
a recomendariam ao produtor; P1 ressalvou que conferiria os valores à mão, e P4 citou a adaptação
da recomendação às fórmulas disponíveis nas revendas da região. Seis responderam que não falta
nada para a sua rotina; as duas exceções são as já citadas: a versão sem conexão (P5) e o
destaque dos campos pendentes (P7).

As notas devem ser lidas como um teto, e não como prova de ausência de problemas. A amostra é
pequena e por conveniência, o questionário foi respondido logo após uma sessão acompanhada pelo
pesquisador, e esses fatores favorecem avaliações positivas. O que o teste sustenta é mais
modesto e mais útil: profissionais do público-alvo, inclusive um sem experiência e um com mais de
dez anos, concluíram sozinhos um cenário de recomendação completo e não encontraram obstáculo que
julgassem mais que estético; e os relatos localizaram um ponto concreto a melhorar, a
sinalização de dados faltantes no formulário.

\subsection{Comparação com ferramenta existente}\label{sec:res_comparacao}
O SIRAS foi posicionado frente ao FertFácil \citep{Tiecher2016, FertFacil2026}, a ferramenta de
recomendação de calagem e adubação mais difundida no RS, adotada pela Emater/RS. A
\autoref{tab:fertfacil} resume a comparação funcional, feita a partir da documentação pública e de
um laudo de exemplo do FertFácil.

\begin{table}[htb]
\centering
\caption{Comparação funcional entre o FertFácil e o SIRAS.}
\label{tab:fertfacil}
\footnotesize
\begin{tabularx}{\columnwidth}{lXX}
\toprule
 & \textbf{FertFácil} & \textbf{SIRAS} \\
\midrule
Base técnica & Manuais RS/SC (2016) e PR (2019) & Manual RS/SC (2016) \\
Escopo & todas as culturas dos dois manuais & 61 culturas e grupos \\
Calagem e adubação & sim & sim \\
Indicação do adubo & uma fórmula, com dose e complemento & fórmula ideal, catálogo e mistura de simples \\
Aptidão edáfica & não & atual e potencial \\
Origem de cada valor & não exibida & regra e tabela do Manual \\
Laudo & sintético, imprimível & essencial e anexo explicativo, imprimível \\
Execução & on-line & local, sem internet \\
Maturidade & em uso desde 2015 & protótipo acadêmico \\
\bottomrule
\end{tabularx}
\end{table}

As duas ferramentas convergem no núcleo --- recomendar calagem e adubação pelo Manual RS/SC --- e
divergem no que entregam em torno dele. O FertFácil cobre mais: dois manuais completos, e não um
recorte de 61 culturas, além de um módulo que compara 11 manuais para a soja. É também uma
ferramenta madura, em uso profissional há mais de uma década. O SIRAS acrescenta três coisas que
o FertFácil não oferece: a estimativa de aptidão edáfica, com o efeito da correção sobre ela; a
indicação da origem de cada número do laudo; e um laudo que, além da recomendação, explica a
interpretação de cada atributo e o caminho até cada dose. O FertFácil entrega a receita; o SIRAS
entrega a receita com a justificativa. As diferenças de escopo e de maturidade são esperadas entre
um produto consolidado e um trabalho de conclusão de curso, e a integração entre recomendação e
aptidão, ausente nas ferramentas analisadas, é a contribuição que este trabalho se propôs a fazer.

A proposta previa também submeter um subconjunto dos casos de validação ao FertFácil e comparar as
saídas, e a mesma restrição impediu uma comparação de usabilidade entre as duas ferramentas. Nenhuma
das duas foi realizada, porque o acesso ao sistema, que exige cadastro, não pôde ser obtido no
período da avaliação. A perda é limitada: as duas ferramentas implementam o mesmo Manual,
e a correção numérica do SIRAS já foi medida contra o oráculo calculado à mão. A comparação
numérica permanece como trabalho futuro, útil sobretudo para localizar pontos de leitura
divergente do Manual entre as duas implementações.

\subsection{Limitações}\label{sec:res_limitacoes}

A aptidão é exclusivamente edáfica: não considera relevo, clima, drenagem nem erosão, e seus
critérios, embora ancorados no Manual e no sistema oficial, foram definidos pelo autor e não
foram validados por especialistas independentes nem em campo. O oráculo das recomendações deriva
do mesmo Manual que o sistema implementa, o que limita a validação à fidelidade aos critérios
publicados. O conjunto de casos é menor que o previsto e desigual entre os grupos. A avaliação de
usabilidade usa amostra pequena e por conveniência, com sessões acompanhadas pelo pesquisador; as
notas de severidade exprimem percepção dos usuários, não gravidade técnica, e a pergunta sobre a
recomendação reuniu o juízo sobre a ferramenta e o juízo sobre a recomendação agronômica. Por fim, o SIRAS é ferramenta de apoio à decisão e não
substitui a responsabilidade técnica de profissional habilitado.

\section{Conclusão}\label{sec:conclusao}

Este trabalho teve por objetivo desenvolver e avaliar o SIRAS, um sistema especialista baseado em
regras que interpreta a análise de solo, recomenda calagem e adubação para 61 culturas e grupos de
culturas segundo o Manual de Calagem e Adubação do RS e de SC e estima a aptidão edáfica da área.
O objetivo foi alcançado: as recomendações concordaram com o oráculo calculado à mão em todos os
25 casos de validação, com limite inferior do intervalo de confiança acima da meta de 80\%, o que
levou à aceitação da hipótese; a implementação da aptidão reproduziu a especificação congelada em
todos os 88 casos de conformidade; e o motor reproduziu as 32 classificações de fósforo e potássio
que o próprio Manual publica no Anexo~1.
Na avaliação de usabilidade, os oito participantes do público-alvo concluíram sozinhos um
cenário de recomendação, nenhum apontou problema além do estético, e todos disseram que usariam a
ferramenta no trabalho.

A principal contribuição é a integração, a partir de uma única análise de solo, de duas
perspectivas que as ferramentas disponíveis tratam em separado: a recomendação de correção e
adubação e a avaliação de aptidão, apresentada antes e depois da correção, de modo que o laudo
mostra não só quanto aplicar, mas o que a aplicação muda. A ela soma-se a rastreabilidade: cada
número do laudo indica a regra e a tabela do Manual que o produziram. A validação mostrou ainda
por que um oráculo independente importa: os casos acrescentados por lacuna de cobertura revelaram
dois defeitos que haviam passado por milhares de testes automatizados.

As limitações mais relevantes são três. A aptidão é exclusivamente edáfica, e seus critérios,
definidos pelo autor a partir do Manual e do sistema oficial de avaliação de aptidão, não foram
validados por especialistas independentes nem em campo. O oráculo das recomendações deriva do
mesmo Manual que o sistema implementa, o que limita a validação à fidelidade aos critérios
publicados. E o conjunto de casos é menor que o previsto e desigual entre os grupos de culturas.
A avaliação de usabilidade, com amostra pequena e por conveniência, indica ausência de obstáculos
graves, mas não permite generalizar a satisfação observada.

Como trabalhos futuros, destacam-se: a validação dos critérios de aptidão por um painel de
agrônomos e por dados de campo; a ampliação do conjunto de casos, sobretudo nos grupos com menos
casos, incluindo laudos reais com recomendação de profissional habilitado; a comparação sistemática
com ferramentas em uso, como o FertFácil, em resultados e em usabilidade; uma versão para celular
que funcione sem conexão, pedida por participante da avaliação; e a inclusão de culturas hoje fora do escopo e de manuais de outras regiões, que a
separação entre base de conhecimento e motor de inferência facilita.


\bibliography{referencias}

% Apêndice: figuras complementares. Não contam no limite de 10 a 15 páginas do corpo.
\appendix
\section{Figuras complementares}\label{sec:apendice_figuras}

As figuras deste apêndice detalham partes do sistema resumidas na \autoref{sec:siras}.

\begin{figure}[htb]
\centering
\figura{F2-arquitetura}{\columnwidth}
\caption{Camadas do SIRAS e dependências entre os pacotes. Nenhuma camada abaixo da interface
depende do Flask.}
\label{fig:arquitetura}
\end{figure}

\begin{figure*}[htb]
\centering
\figura{F6-interpretacao}{\textwidth}
\caption{Interpretação de fósforo pela classe de teor de argila e de potássio pela CTC a pH
7,0, para o grupo de exigência das culturas de grãos (Manual 2016, Tab. 6.4 e 6.9). O ponto
marca a Gleba~4 do Anexo~1.}
\label{fig:interpretacao}
\end{figure*}

\begin{figure*}[htb]
\centering
\figura{F7-adubacao}{\textwidth}
\caption{Os dois caminhos da adubação: o algoritmo dos grãos e a dose por classe dos demais
grupos, parametrizada pelas variáveis de cada um.}
\label{fig:adubacao}
\end{figure*}

\begin{figure*}[htb]
\centering
\figura{F8-formulacao}{\textwidth}
\caption{Os três caminhos da formulação do adubo e a escolha da dose pelos pontos de quebra,
ilustrados com a saída do motor para a Gleba~4.}
\label{fig:formulacao}
\end{figure*}

\begin{figure}[htb]
\centering
\figura{F10-trilha}{\columnwidth}
\caption{Anatomia de um passo da trilha de inferência e sua apresentação no laudo.}
\label{fig:trilha}
\end{figure}

\begin{figure*}[htb]
\centering
\figura{F12-fluxo-de-uso}{\textwidth}
\caption{Fluxo de uso do SIRAS na perspectiva do usuário.}
\label{fig:fluxo}
\end{figure*}

\begin{figure}[htb]
\centering
\figura{F14-usabilidade}{\columnwidth}
\caption{Desenho da sessão de avaliação de usabilidade.}
\label{fig:sessao}
\end{figure}


\end{document}
